\documentclass[11pt]{article}

\usepackage[margin=1in]{geometry}
\usepackage{amsmath,amssymb,amsthm}
\usepackage{graphicx}
\usepackage[hidelinks]{hyperref}

\newtheorem{theorem}{Theorem}
\newtheorem{lemma}{Lemma}
\newtheorem{remark}{Remark}

\DeclareMathOperator{\Lip}{Lip}
\DeclareMathOperator{\supp}{supp}
\newcommand{\FM}{\mathcal F(M)}
\newcommand{\MU}{M^{\mathcal U}}
\newcommand{\SM}{\mathcal S}
\newcommand{\N}{\mathbb N}
\newcommand{\ip}[2]{\langle #1,#2\rangle}

\title{A nonseparable extended support for a Lipschitz functional}
\author{Candidate counterexample to a subquestion in arXiv:2009.07663}
\date{2026-06-21}

\begin{document}
\maketitle

\section*{Status}

This packet records a candidate counterexample, likely valid, to the
separability part of Question 1 in Aliaga--Perneck\'a \cite{AP}.  It does not
disprove the weaker countable-chain-condition alternative in that question.
Indeed, the support constructed below still satisfies ccc.

\section*{Source question}

The source paper is Ramon J. Aliaga and Eva Perneck\'a, ``Integral
representation and supports of functionals on Lipschitz spaces'',
arXiv:2009.07663 \cite{AP}.  In the concluding remarks, after proving that
\(\SM(\phi)\cap M\) is separable for every
\(\phi\in \Lip_0(M)^*\), they ask:

\begin{center}
\includegraphics[width=0.94\linewidth]{figures/open_problem_crop.png}
\end{center}

We answer the first half of Question 1 negatively: \(\SM(\phi)\) need not be a
separable subset of \(\MU\).

\section*{Definitions used}

For a pointed complete metric space \(M\), \(\Lip_0(M)\) denotes the real
Banach space of Lipschitz functions \(f:M\to\mathbb R\) with \(f(0)=0\), with
norm the Lipschitz constant.  Its canonical predual is the Lipschitz-free space
\(\FM\), and \(\Lip_0(M)^*\) is the bidual \(\mathcal F(M)^{**}\).

Aliaga--Perneck\'a define the extended support of
\(\phi\in\Lip_0(M)^*\) as
\[
\SM(\phi)=
\bigcap\{K\subset \MU: K\text{ compact and }\phi\in\mathcal I(K)^\perp\},
\]
where \(\MU\) is the uniform compactification and
\[
\mathcal I(K)=\{f\in\Lip_0(M): f^{\mathcal U}=0\text{ on }K\}.
\]
We use their Proposition 4.7 \cite[Proposition 4.7]{AP}: if
\(\phi\in\Lip_0(M)^*\) is induced by an almost Radon measure \(\mu\) on
\(\MU\), then
\[
        \SM(\phi)=\supp(\mu).
\]

\section*{Idea of the proof}

Take a nonseparable compact space that nevertheless carries a full-support
Radon probability measure: the Cantor cube \(K=\{0,1\}^I\) for uncountable
\(I\), with its Bernoulli product measure.  Let \(\Gamma\) be a dense subset of
\(K\), but forget the topology inherited from \(K\) and give \(\Gamma\) the
bounded discrete metric.  Then the uniform compactification of \(\Gamma\) is
\(\beta\Gamma\), because every bounded function on a uniformly discrete space
is uniformly continuous.

The dense inclusion \(\Gamma\to K\) is uniformly continuous from the discrete
metric and extends to a continuous surjection
\(\pi:\beta\Gamma\to K\).  Lift the Bernoulli measure from \(K\) to a Radon
probability \(\nu\) on \(\beta\Gamma\).  Its support cannot be separable:
otherwise its continuous image under \(\pi\) would be a separable compact
subset of \(K\), but the image is all of \(K\).

Finally let \(\phi\) be the functional on \(\Lip_0(\Gamma)\) obtained by
integrating Stone--Cech extensions against \(\nu\).  Proposition 4.7 of the
source paper identifies \(\SM(\phi)\) with \(\supp(\nu)\), so the extended
support is nonseparable.

\begin{lemma}\label{lem:cube}
Let \(I\) be an uncountable set and \(K=\{0,1\}^I\) with the product topology.
Then \(K\) is nonseparable and carries a Radon probability measure with full
support.
\end{lemma}

\begin{proof}
Let \(\lambda=(\frac12\delta_0+\frac12\delta_1)^I\) be the product probability
measure.  Every nonempty basic open cylinder has measure \(2^{-n}\) for some
finite \(n\), so \(\lambda\) has full support.  Since \(K\) is compact
Hausdorff, \(\lambda\) is Radon.

For the nonseparability, we use the standard density calculation.  If
\(D\subset 2^I\) is dense, then different coordinates must have different
restrictions to \(D\).  Indeed, if coordinates \(i\ne j\) had
\(x(i)=x(j)\) for every \(x\in D\), then \(D\) would miss the nonempty basic
open set \(\{x:x(i)=0,\ x(j)=1\}\).  Thus the map
\[
        I\longrightarrow 2^D,\qquad i\longmapsto (x(i))_{x\in D}
\]
is injective, so \(|I|\leq 2^{|D|}\).  In this packet we fix
\(|I|>2^{\aleph_0}\), and then no countable \(D\) can be dense.
\end{proof}

\begin{lemma}\label{lem:lift}
Let \(\Gamma\) be a dense subset of the compact space \(K\), equipped with the
discrete metric \(d(\gamma,\eta)=1\) for \(\gamma\ne\eta\).  If
\(\lambda\) is a Radon probability measure on \(K\), then there is a Radon
probability measure \(\nu\) on \(\beta\Gamma\) such that
\(\pi_*\nu=\lambda\), where
\(\pi:\beta\Gamma\to K\) is the continuous extension of the inclusion
\(\Gamma\hookrightarrow K\).
\end{lemma}

\begin{proof}
The inclusion \(\Gamma\to K\) is uniformly continuous because \(\Gamma\) has
the discrete metric.  Hence it extends continuously to the Stone--Cech
compactification \(\beta\Gamma\).  Its image is compact and contains the dense
set \(\Gamma\), so it is all of \(K\).

Define a positive norm-one functional on the subspace
\(\{g\circ\pi:g\in C(K)\}\subset C(\beta\Gamma)\) by
\[
        g\circ\pi\longmapsto \int_K g\,d\lambda .
\]
By the positive Hahn--Banach extension theorem, it extends to a positive
norm-one functional on \(C(\beta\Gamma)\).  The Riesz representation theorem
gives a Radon probability \(\nu\) on \(\beta\Gamma\).  For every
\(g\in C(K)\),
\[
        \int_{\beta\Gamma} g\circ\pi\,d\nu=\int_K g\,d\lambda,
\]
which is precisely \(\pi_*\nu=\lambda\).
\end{proof}

\begin{theorem}\label{thm:counterexample}
There exist a bounded complete metric space \(M\) and a positive functional
\(\phi\in\Lip_0(M)^*\) such that \(\SM(\phi)\) is not separable.
\end{theorem}

\begin{proof}
Choose a set \(I\) with \(|I|>2^{\aleph_0}\), put \(K=2^I\), and let
\(\lambda\) be the full-support Bernoulli product measure from
Lemma~\ref{lem:cube}.  Let \(\Gamma\subset K\) be the set of finite-support
points:
\[
        \Gamma=\{x\in 2^I:\{i:x(i)=1\}\text{ is finite}\}.
\]
This set is dense in \(K\).  Choose the zero function in \(\Gamma\) as base
point and equip \(M=\Gamma\) with the discrete metric
\[
        d(x,y)=\begin{cases}0,&x=y,\\1,&x\ne y.\end{cases}
\]
Then \(M\) is bounded and complete.  Since \(M\) is uniformly discrete, every
bounded function on \(M\) is uniformly continuous, so its uniform
compactification is \(\MU=\beta M\).  Also \(\Lip_0(M)\) consists exactly of
the bounded real functions on \(M\) vanishing at the base point, with an
equivalent norm.

Let \(\pi:\beta M\to K\) be the extension of the inclusion, and let \(\nu\) be
a Radon probability on \(\beta M\) with \(\pi_*\nu=\lambda\), supplied by
Lemma~\ref{lem:lift}.  Define
\[
        \ip{f}{\phi}=\int_{\beta M} f^{\mathcal U}\,d\nu,
        \qquad f\in\Lip_0(M).
\]
This is a positive bounded functional on \(\Lip_0(M)\), induced by the Radon
measure \(\nu\) on \(\MU\).  By \cite[Proposition 4.7]{AP},
\[
        \SM(\phi)=\supp(\nu).
\]

We claim that \(\supp(\nu)\) is nonseparable.  Since \(\pi_*\nu=\lambda\) and
\(\supp(\lambda)=K\), one has
\[
        K=\supp(\lambda)\subseteq \pi(\supp(\nu)).
\]
The reverse inclusion is automatic because \(\pi(\supp(\nu))\subset K\), hence
\(\pi(\supp(\nu))=K\).  If \(\supp(\nu)\) were separable, then its continuous
image \(K\) would be separable, contradicting Lemma~\ref{lem:cube}.  Thus
\(\SM(\phi)\) is nonseparable.
\end{proof}

\begin{remark}[What remains open]
This construction does not refute the ccc alternative in Question 1.  In fact,
\(\SM(\phi)=\supp(\nu)\) has ccc: if \(\{U_i\}_{i\in A}\) is a pairwise
disjoint family of nonempty relatively open subsets of \(\supp(\nu)\), then
\(\nu(U_i)>0\) for every \(i\), and finite additivity over finite subfamilies
implies that only countably many such \(U_i\)'s can exist.
\end{remark}

\section*{Verification and novelty check}

The proof is purely structural.  The local run indexes were searched for
arXiv id \(2009.07663\), the title, ``extended support'', ``countable chain
condition'', ``Cantor cube'', and ``nonseparable support''; no duplicate packet
or attempt was found.  Bounded web searches on 2026-06-21 for exact phrases
from Question 1 and for the paper title found no later explicit solution of
the separability subquestion.

Human review should focus on the measure-lifting step in Lemma~\ref{lem:lift}
and on the invocation of Aliaga--Perneck\'a Proposition 4.7 identifying
extended support with the support of the inducing Radon measure.

\begin{thebibliography}{9}
\bibitem{AP}
R. J. Aliaga and E. Perneck\'a,
\emph{Integral representation and supports of functionals on Lipschitz spaces},
arXiv:2009.07663, 2020.
\end{thebibliography}

\end{document}
